%-------------------------
% One-page CV -- Italian. Layout and shared macros live in preamble.tex.
% Keep resume.tex in sync with every content change made here.
%-------------------------
\documentclass[a4paper,11pt,italian]{article}
%-------------------------
% Resume in Latex
% Author : Sourabh Bajaj
% Website: https://github.com/sb2nov/resume
% License : MIT
%------------------------
%
% Shared preamble for every variant of the one-page CV.
%
%   resume.tex            English
%   resume-it.tex         Italian
%   resume-signed.tex     English + GDPR clause and signature
%   resume-it-signed.tex  Italian + GDPR clause and signature
%
% The document language comes from the \documentclass option
% (english / italian) and is picked up by babel. A signed variant
% defines \cvsigned before loading the CV. Every variant must fit
% on ONE page.
%------------------------


\usepackage{babel}
\usepackage{latexsym}
\usepackage{fontawesome}
\usepackage{graphicx}
\usepackage[empty]{fullpage}
\usepackage{titlesec}
\usepackage{marvosym}
\usepackage[usenames,dvipsnames]{color}
\usepackage{verbatim}
\usepackage{enumitem}
\usepackage[hidelinks]{hyperref}
\usepackage{fancyhdr}
\usepackage{comment}
\usepackage{fontspec}
\usepackage{setspace}
\usepackage{xcolor}

\DeclareRobustCommand{\faIcon}[1]{\faicon{#1}}

\definecolor{navy}{rgb}{0.0, 0.0, 0.6}
\definecolor{blue}{rgb}{0.0, 0.0, 0.8}
\setstretch{1.0}

\pagestyle{fancy}
\fancyhf{} % clear all header and footer fields
\fancyfoot{}
\renewcommand{\headrulewidth}{0pt}
\renewcommand{\footrulewidth}{0pt}

% Adjust margins
\addtolength{\oddsidemargin}{-.6in}
\addtolength{\evensidemargin}{-.6in}
\addtolength{\textwidth}{1.25in}
\addtolength{\topmargin}{-.6in}
\addtolength{\textheight}{1.2in}

\urlstyle{same}

\raggedbottom
\raggedright
\setlength{\tabcolsep}{0in}

% Sections formatting
\titleformat{\section}{
  \vspace{-4pt}\scshape\raggedright\large
}{}{0em}{}[\color{black}\titlerule \vspace{3pt}]
\titlespacing*{\section}{0pt}{4pt}{4pt}

%-------------------------
% Custom commands
\newcommand{\resumeItem}[2]{
  \item\small{
    \textbf{#1}{: #2 \vspace{-2pt}}
  }
}

\newcommand{\resumeSubheading}[4]{
  \vspace{1pt}\item[]
    \begin{tabular*}{\linewidth}{@{}l@{\extracolsep{\fill}}r@{}}
      \textbf{#1} & #2 \\
      \textit{#3} & \textit{#4} \\
    \end{tabular*}\vspace{-4pt}
}

\newcommand{\resumeSubItem}[2]{\resumeItem{#1}{#2}\vspace{-3pt}}

\renewcommand{\labelitemii}{$\circ$}

\newcommand{\resumeSubHeadingListStart}{\begin{itemize}[leftmargin=*,label={}]}
\newcommand{\resumeSubHeadingListEnd}{\end{itemize}\vspace{-3pt}}
\newcommand{\resumeItemListStart}{\begin{itemize}[leftmargin=12pt,label={$\bullet$}]}
\newcommand{\resumeItemListEnd}{\end{itemize}\vspace{-3pt}}
%-------------------------------------------
%%%%%%%%%%%%  CV STARTS HERE  %%%%%%%%%%%%%%

\hypersetup{
  pdfauthor={Alessandro Monticelli},
  pdfsubject={Curriculum Vitae},
  pdfkeywords={Alessandro Monticelli, CV, Software Engineer, Embedded Systems, LoRa, RAG},
}

% The page is full: keep the signature block compact.
\def\cvsignatureskip{2pt}
\def\cvsignatureheight{0.8cm}
%-------------------------
% Signed variants: GDPR consent clause, place/date and signature.
%
% \cvsignatureblock prints nothing unless \cvsigned is defined
% (see resume-signed.tex / resume-it-signed.tex).
%
% The signature image is read from signature.png when it exists.
% In CI (.github/workflows/signed.yml) it is decrypted from the committed
% signature.png.gpg with the SIGNATURE_PASSPHRASE secret, so the plain
% scan is never committed. Without it a blank signature line is printed.
%
% Set the place shown before the date with \def\cvplace{City}
% (left empty, only the date is printed).
%-------------------------

\ifdefined\cvplace\else\def\cvplace{}\fi
% Vertical space before the block and height of the signature image;
% a preamble can \def these before %-------------------------
% Signed variants: GDPR consent clause, place/date and signature.
%
% \cvsignatureblock prints nothing unless \cvsigned is defined
% (see resume-signed.tex / resume-it-signed.tex).
%
% The signature image is read from signature.png when it exists.
% In CI (.github/workflows/signed.yml) it is decrypted from the committed
% signature.png.gpg with the SIGNATURE_PASSPHRASE secret, so the plain
% scan is never committed. Without it a blank signature line is printed.
%
% Set the place shown before the date with \def\cvplace{City}
% (left empty, only the date is printed).
%-------------------------

\ifdefined\cvplace\else\def\cvplace{}\fi
% Vertical space before the block and height of the signature image;
% a preamble can \def these before \input{signature} to save space.
\ifdefined\cvsignatureskip\else\def\cvsignatureskip{8pt}\fi
\ifdefined\cvsignatureheight\else\def\cvsignatureheight{1cm}\fi

\newcommand{\cvconsenttext}{%
  \iflanguage{italian}{%
    Autorizzo il trattamento dei dati personali contenuti nel presente CV
    ai sensi del D.Lgs.~196/2003 e del Regolamento (UE) 2016/679 (GDPR).%
  }{%
    I authorize the processing of the personal data contained in this CV
    pursuant to Legislative Decree 196/2003 and EU Regulation 2016/679 (GDPR).%
  }%
}

\newcommand{\cvsignatureblock}{%
  \ifdefined\cvsigned
    \par\vspace{\cvsignatureskip}%
    \noindent\begin{minipage}[b]{0.66\linewidth}%
      \footnotesize\raggedright\cvconsenttext\par\vspace{2pt}%
      \ifx\cvplace\empty\else\cvplace, \fi\today
    \end{minipage}\hfill
    \begin{minipage}[b]{0.3\linewidth}%
      \centering
      \IfFileExists{signature.png}%
        {\includegraphics[height=\cvsignatureheight,width=\linewidth,keepaspectratio]{signature.png}}%
        {\rule{0pt}{\cvsignatureheight}}\\[-2pt]
      \rule{\linewidth}{0.4pt}\\
      \footnotesize Alessandro Monticelli
    \end{minipage}%
  \fi
}
 to save space.
\ifdefined\cvsignatureskip\else\def\cvsignatureskip{8pt}\fi
\ifdefined\cvsignatureheight\else\def\cvsignatureheight{1cm}\fi

\newcommand{\cvconsenttext}{%
  \iflanguage{italian}{%
    Autorizzo il trattamento dei dati personali contenuti nel presente CV
    ai sensi del D.Lgs.~196/2003 e del Regolamento (UE) 2016/679 (GDPR).%
  }{%
    I authorize the processing of the personal data contained in this CV
    pursuant to Legislative Decree 196/2003 and EU Regulation 2016/679 (GDPR).%
  }%
}

\newcommand{\cvsignatureblock}{%
  \ifdefined\cvsigned
    \par\vspace{\cvsignatureskip}%
    \noindent\begin{minipage}[b]{0.66\linewidth}%
      \footnotesize\raggedright\cvconsenttext\par\vspace{2pt}%
      \ifx\cvplace\empty\else\cvplace, \fi\today
    \end{minipage}\hfill
    \begin{minipage}[b]{0.3\linewidth}%
      \centering
      \IfFileExists{signature.png}%
        {\includegraphics[height=\cvsignatureheight,width=\linewidth,keepaspectratio]{signature.png}}%
        {\rule{0pt}{\cvsignatureheight}}\\[-2pt]
      \rule{\linewidth}{0.4pt}\\
      \footnotesize Alessandro Monticelli
    \end{minipage}%
  \fi
}


\hypersetup{pdftitle={Alessandro Monticelli -- Curriculum Vitae (IT)}, pdflang={it-IT}}

\begin{document}
\setlength{\footskip}{5pt}
%----------HEADING-----------------
\begin{center}
  \textbf{\Large Alessandro Monticelli} \\ \vspace{3pt}
  {\small{\faIcon{phone}} +39-340-053-1380 \quad
    \href{mailto:alessandro.monticelli@live.it}{{\faIcon{envelope}} \underline{alessandro.monticelli@live.it}} \quad
    \href{https://www.linkedin.com/in/ale-mont}{{\faIcon{linkedin}} \underline{Alessandro Monticelli}} \quad
    \href{https://github.com/aleemont1}{{\faIcon{github}} \underline{aleemont1}}}
\end{center}
\vspace{-6pt}

%-----------EXPERIENCE-----------------
\section{Esperienze professionali}
\resumeSubHeadingListStart
    \resumeSubheading{Sviluppatore Software}{Faenza, Italia}{S.EL.E.D. Soc. Coop.}{Settembre 2026 -- In corso}
    \resumeItemListStart
        \item\small{Realizzato in SELED \textbf{RagCore}, motore RAG generico per contratti e referti di laboratorio (hit@1 0,948), configurato tramite Knowledge Base dichiarative e isolato dal database.}
        \item\small{Ridotto di 15 volte i costi API di acquisizione documenti nelle pipeline ETL di Power Automate con Claude Haiku 5.5, prompt caching e output strutturati via tool calling (al posto di lunghi prompt con schemi JSON).}
        \item\small{Automatizzato il monitoraggio delle disdette contrattuali, eliminando il rischio di penali da rinnovi mancati per Caviro; rifattorizzato la PowerApp per gli XML delle fatture estere verso lo SdI (LINKER) con aggiornamenti IVA massivi.}
    \resumeItemListEnd

    \resumeSubheading{Software Engineer freelance}{Da remoto}{Libero professionista}{Giugno 2026 -- In corso}
    \resumeItemListStart
        \item\small{Sviluppato \href{https://github.com/aleemont1/supermarket-scraper}{GDOScraper}, pipeline ETL modulare che estrae e normalizza le promozioni della GDO per la startup italiana Foodify, tramite scraping dinamico di REST API, parsing spaziale di volantini PDF e multiprocessing.}
    \resumeItemListEnd

    \resumeSubheading{Studente ricercatore}{Cesena, Italia}{Università di Bologna (supervisione del prof. Andrea Piroddi)}{Luglio 2025 -- In corso}
    \resumeItemListStart
        \item\small{Progettato \href{https://github.com/aleemont1/LAMP}{LAMP (LoRaMultiPacket)}, protocollo di trasporto affidabile in C++17 con SACK-ARQ a singola bitmap (fino a 62,2~kB); valutato su 295.800 simulazioni ns-3 (+22--35\% di portata al 90\% di PDR in ambito urbano) e su banco hardware ESP32-S3/SX1262. Articolo IEEE in stesura.}
    \resumeItemListEnd

    \resumeSubheading{Project Manager}{Forlì, Italia}{Aurora Rocketry Team - Unibo}{Aprile 2024 -- In corso}
    \resumeItemListStart
        \item\small{Coordino un team di 7 sviluppatori sul flight computer ESP32 di \href{https://github.com/AuroraRocketryTeam/atlas}{Atlas / Nemesis} (razzo sonda), con telemetria LoRa\textsuperscript{\textcopyright} in tempo reale, sensor fusion e rilevamento dell'apogeo via EKF (base della tesi triennale).}
    \resumeItemListEnd
\resumeSubHeadingListEnd

%-----------EDUCATION-----------------
\section{Istruzione}
\resumeSubHeadingListStart
    \resumeSubheading{``Alma Mater Studiorum'' - Università di Bologna}{Cesena, Italia}
    {Laurea Magistrale in Ingegneria e Scienze Informatiche}{Set 2025 -- In corso}
    \resumeSubheading{Reykjavik University (H\'ask\'olinn \'i Reykjav\'ik)}{Reykjavík, Islanda}
    {Erasmus+ per la Laurea Magistrale in Computer Science}{Gen 2026 -- Mag 2026}
    \resumeItemListStart
        \resumeItem{Insegnamenti rilevanti}{Distributed and Concurrent Programming (Go), Cybersecurity Management and Compliance (ISO27001, ISO42001), Deep Learning, Big Data Management}
    \resumeItemListEnd
    \resumeSubheading{``Alma Mater Studiorum'' - Università di Bologna}{Cesena, Italia}
    {Laurea Triennale in Ingegneria e Scienze Informatiche}{Set 2021 -- Lug 2025}
    \resumeItemListStart
        \resumeItem{Insegnamenti rilevanti}{Sistemi Embedded e IoT, Programmazione di Reti, Reti di Telecomunicazione, High Performance Computing, Programmazione ad Oggetti}
    \resumeItemListEnd
\resumeSubHeadingListEnd

%-----------SKILLS-----------------
\section{Competenze}
\resumeSubHeadingListStart
    \resumeSubItem{Programmazione}{C++, C, Python, Java, C\# (.NET 8.0), Go (base), SQL.}
    \resumeSubItem{AI \& Data Engineering}{Architetture RAG, PostgreSQL (\texttt{pgvector}, \texttt{tsvector}), scikit-learn, API di LLM (Prompt Caching), FastAPI.}
    \resumeSubItem{Embedded \& IoT}{ESP32 (FreeRTOS), Arduino, LoRa\textsuperscript{\textcopyright}, MQTT, Sensor Fusion, scheduling RTOS.}
    \resumeSubItem{Systems \& Cloud}{Docker, Kubernetes (base), CI/CD (GitHub Actions), Power Platform, Linux (Desktop \& Server), Windows Server.}
\resumeSubHeadingListEnd

%-----------PUBLICATIONS & PROJECTS-----------------
\section{Progetti \& Pubblicazioni}
\resumeSubHeadingListStart
    \resumeSubItem{\href{https://github.com/aleemont1/lora-paper}{Articolo di ricerca su LoRa (in preparazione)}}{Monticelli A., Piroddi A., \textit{“LAMP: A Selective-ACK Transport Protocol for Reliable Bulk Data Transmission over LoRa Networks”}, in preparazione per la sottomissione a IEEE, 2026.}
    \resumeSubItem{\href{https://aleemont1.github.io/bachelor-thesis/pdf/tesi.pdf}{Tesi di Laurea Triennale}}{\textit{“Sviluppo e ottimizzazione di un sistema di telemetria per razzi -- Architettura e trasmissione dati tramite LoRa\textsuperscript{\textcopyright}”}, Università di Bologna, 2025.}
    \resumeSubItem{\href{https://github.com/AuroraRocketryTeam/atlas}{Nemesis (Flight Computer)}}
    {Software di volo ESP32 per razzi ad alta potenza, con telemetria LoRa\textsuperscript{\textcopyright} in tempo reale, sensor fusion e rilevamento dell'apogeo via EKF.}
\resumeSubHeadingListEnd

\cvsignatureblock
\end{document}
